\section{Transfer of the conclusions to characteristic-zero fields}
\label{bl:sec:fields}

The remaining task is to recover generation of an actual line bundle
on the original variety over an arbitrary algebraically closed field
of characteristic zero. We descend the specified divisor data to a
countable algebraically closed field, embed that field into $\C$,
and detect generation by the evaluation cokernel. Faithful flatness
then transfers the same Cartier multiple back. The auxiliary minimal
models used over $\C$ play no role in this descent.

\begin{lemma}[Transfer of specified divisor data]
\label[lemma]{bl:lem:field-transfer}
Let $\kk$ be an algebraically closed field of characteristic zero.
A finite collection of projective varieties over $\kk$, rational
divisors, specified Cartier multiples, morphisms, line bundles,
sections, and identities between these data can be defined over
a countable algebraically closed subfield $\kk_0\subset\kk$.
The field $\kk_0$ admits an embedding into $\C$.

For the data considered here, one can include a log resolution and
its discrepancy identities in this descent.  Log canonicity, and
klt singularities outside a specified reduced boundary, then hold
after extension from $\kk_0$ to either $\kk$ or $\C$.  For a
rational Cartier divisor on a projective variety over $\kk_0$,
nefness, Kodaira dimension, and generation of any fixed Cartier
multiple are unchanged under algebraically closed field extension.
\end{lemma}

\begin{proof}
Choose finitely many equations and transition functions defining
the schemes, morphisms, line bundles, and sections in question.
Equalities and isomorphisms among them also have finite
presentation data.  They are defined over a finitely generated
subfield $\kk_1\subset\kk$; enlarge it finitely whenever another
one of the specified data is added.  Take for $\kk_0$ the algebraic
closure of $\kk_1$ inside $\kk$.  It is algebraically closed and
countable.  Choosing images in $\C$ of a transcendence basis of
$\kk_1/\Q$, and then extending over its algebraic closure, gives
an embedding $\kk_0\hookrightarrow\C$.

Include the individual prime divisors, a projective log resolution,
its exceptional divisors, and all the required Cartier powers in
this construction.  The resulting varieties over $\kk_0$ are
normal or smooth as appropriate; these properties descend along
the faithfully flat extension to $\kk$.  Over the algebraically
closed field $\kk_0$, normal varieties are geometrically normal
(see \cite[Tag~0C3M]{StacksProject}), and integral varieties are
geometrically integral.  The divisors, simple normal crossing
conditions on the resolution, and birationality of the chosen
morphisms consequently persist under extension.  One can also
include an open set on which a birational morphism is an
isomorphism.

For canonical divisors, choose compatible rational top forms on
the smooth loci and on the resolution, and descend them together
with the rational functions specifying the divisor comparisons.
Pullback identities for the selected rational Cartier divisors
then remain identities after extension.  In particular a formula
for the log discrepancies on the chosen resolution is unchanged.
The finitely many coefficient inequalities on its snc boundary
test lc singularities.  On the inverse image of the complement
of the specified floor, the strict inequalities test klt
singularities.  They therefore give the asserted singularity
properties on the complex extension as well.  Only the specified
Cartier powers are being transferred; this argument asserts no
general preservation of $\Q$-factoriality.

We give the nefness argument, since it involves all curves and
is not itself a finite list of witnesses.  Let $X_0$ be projective
over $\kk_0$, let $L_0$ be a line bundle, and let $\kk'$ be an
algebraically closed extension.  An integral curve over $\kk_0$
remains integral after extension and retains its $L_0$-degree.
Thus nefness over $\kk'$ implies nefness over $\kk_0$.
Conversely, suppose an integral curve $C\subset X_{\kk'}$ has
negative degree.  It is defined over a finitely generated field
extension of $\kk_0$.  Spread it as a closed subscheme of
$X_0\times T$, where $T$ is an integral finite-type
$\kk_0$-scheme with that function field.  After shrinking $T$,
the family is flat with geometrically integral one-dimensional
fibers, and the degree of $L_0$ on the fibers is constant.  These
assertions follow from geometric integrality of the generic
fiber, generic flatness, and constancy of the relevant Hilbert
polynomials.  Every nonempty such $T$ has a closed
$\kk_0$-point.  Its fiber is therefore a negative curve on
$X_0$, a contradiction.  Apply this argument to a Cartier
multiple for rational divisors.

For every integer $m$ making $mD_0$ Cartier, flat base change
gives
\begin{equation}
 H^0(X_0,\OO_{X_0}(mD_0))\otimes_{\kk_0}\kk'
   \simeq H^0(X_{\kk'},\OO_{X_{\kk'}}(mD_{\kk'}));
 \label{bl:eq:fields-sections}
\end{equation}
see \cite[Tag~02KH]{StacksProject}.  The complete linear-system
maps therefore base change, as do their image closures, so their
image dimensions and the Kodaira dimension are unchanged.
Alternatively, their spaces of sections have the same dimensions
in every divisible degree.  The evaluation map for $mD_0$ also
base changes to the evaluation map for $mD_{\kk'}$.  Its
cokernel vanishes exactly when its faithfully flat pullback
vanishes.  This proves generation in a fixed degree, and thus
semiampleness, in both directions.  The same argument applies to
line bundles on the whole reduced boundary scheme, even when
that scheme is reducible.
\end{proof}

\begin{proof}[Completion of Theorem~\ref{bl:thm:supported}]
Start with its given data over $\kk$.  In
Lemma~\ref{bl:lem:field-transfer}, include the pair, every component
of its boundary and of $G$, the exact identity
$G\sim q(K_V+C)$, and the section defining the nonzero effective
Cartier divisor $G$.  Include also a generated positive power of
$\OO_V(G)|_S$ with finitely many generating sections on the
entire $S=G_{\mathrm{red}}$.  Reduced supports and their
inclusions are part of these data, so $G$ stays nonzero and
$\Supp G\subset\Supp\lfloor C\rfloor$ after extension.

Because the original variety is $\Q$-factorial, $K_V$ and each
of the finitely many boundary components have individual Cartier
multiples; include these witnesses.  The descended discrepancy
data give lc singularities and klt singularities outside the
floor over $\C$.  These are exactly the weaker complex
hypotheses in Proposition~\ref{bl:prop:complex-supported}; global
$\Q$-factoriality of the complex extension is unnecessary.
That proposition gives $\kappa(K_V+C)>0$ on the complex
extension.  Equation~\eqref{bl:eq:fields-sections} transfers this
conclusion first to $\kk_0$ and then to $\kk$, proving the
supported theorem as stated.
\end{proof}

\begin{corollary}[Abundance after nonvanishing over characteristic-zero fields]
\label{cor:abundance-after-nonvanishing-fields}
Let $(X,\Delta)$ be a projective lc rational pair of dimension at most
four over an algebraically closed field $\kk$ of characteristic zero.
Suppose that $D=K_X+\Delta$ is $\Q$-Cartier and nef and
$\kappa(X,D)\geq0$. Then $D$ is semiample, and if $\kappa(X,D)=0$,
then $D\sim_{\Q}0$.
\end{corollary}

\begin{proof}[Proof of Corollary~\ref{cor:abundance-after-nonvanishing-fields} and completion of Theorem~\ref{bl:thm:conditional}]
For the four-dimensional assertion, descend the original lc pair
and its rational Cartier adjoint $D$, with a log resolution and
its discrepancy identities.  No separate Cartier property for
$K_X$ or the boundary components is needed for this application.
The complex extension is lc, $D$ is nef, and its Kodaira
dimension is the same.  The complex case proved in
Section~\ref{bl:sec:abundance} makes one positive Cartier multiple
of $D$ generated.  Generation in this particular degree descends
to $\kk_0$ and then extends to $\kk$.  If $\kappa(D)=0$, the
generated line bundle has a one-dimensional space of sections;
its nonzero generating section is nowhere vanishing.  Hence it
is trivial, and $D\sim_{\Q}0$ over $\kk$.

For the conditional assertion, keep its smooth canonical
nonvanishing premise over $\C$ exactly as stated.  Descend a
given nef projective lc rational pair over $\kk$, extend to
$\C$, and apply the conditional complex result there.  Its
proof may use very general complex fibers and complex minimal
models.  Only the resulting generation of a fixed positive
multiple of the original adjoint is transferred back by
Lemma~\ref{bl:lem:field-transfer}.  This needs neither very general
points over the countable field $\kk_0$ nor descent of the
auxiliary minimal models, and completes both theorems.
\end{proof}


\begin{proof}[Proof of \Cref{thm:full-abundance}]
First work over $\C$. Fix a Cartier index $r>0$ for the given
nef adjoint $D$. The exact original-variety lc nonvanishing result
\cite[Corollary~1.2]{OpenAINonvanishing2026}
produces a nonzero section of $mrD$ for some $m>0$. Thus
$\kappa(X,D)\geq0$, and Theorem~\ref{thm:abundance-after-nonvanishing}
proves semiampleness, with $\Q$-linear triviality when $\kappa=0$.
The cited nonvanishing corollary uses smooth fourfold nonvanishing,
Hashizume's reduction, and the independent finite rational-feasibility
lemma \cite[Lemma~8.1]{OpenAINonvanishing2026}.
It supplies a section on the original normal variety and does not use
any supported lifting or abundance theorem from this paper.

Over an arbitrary algebraically closed characteristic-zero field,
descend the original pair, a specified Cartier multiple of its
adjoint, and a log resolution with its discrepancy identities as in
\Cref{bl:lem:field-transfer}. The complex extension is lc and its
adjoint is nef. The complex conclusion makes one positive Cartier
multiple generated. The evaluation-cokernel argument of that lemma
transfers this precise generation to the original field. Kodaira
dimension transfers as well. If it is zero, the generated multiple
has a one-dimensional space of sections; its nonzero generating
section is nowhere vanishing, hence trivializes that line bundle.
This proves the stated $\Q$-linear triviality on the original variety.
\end{proof}
